\section{总结与延伸}

\subsection{一条贯穿全课的推理链}

本课可以压缩为六步。第一，为参数、optimizer state、gradient、activation 与数据流分别记账。第二，沿最匹配的逻辑轴做 sharding。第三，写出恢复单设备语义所需的 collective 或 point-to-point 依赖。第四，判断通信位于 critical path 还是可 overlap。第五，把逻辑 axes 映射到真实 NVLink、NVSwitch、InfiniBand、PCIe 与 storage 拓扑。第六，用 profiler、吞吐、峰值内存、网络利用率和收敛曲线验证，而不是凭配置名称验收。

\begin{importantbox}{Ultra-scale 的核心不是更多 GPU，而是更少无效等待}
增加设备只增加潜在并行度。只有当状态被正确分片、collective 被正确分组、micro-batch 能填满 schedule、路由负载平衡、数据与 checkpoint 不饿死计算时，潜在并行度才会变成有效训练吞吐。
\end{importantbox}

\subsection{实践检查清单}

\begin{enumerate}
  \item 记录每类状态的 dtype、shape、replication factor 与生命周期。
  \item 为每个 parallel axis 标出 process group、物理链路与关键 collective。
  \item 分别测量 compute-only、communication-only 与 end-to-end step time。
  \item 检查 overlap 是否降低总时间，而不是只让时间线重叠。
  \item 记录 straggler、router imbalance、pipeline bubble 和 input/checkpoint stall。
  \item 用最小规模复现单卡 loss，再逐轴增加并行并比较数值与吞吐。
  \item 保存可恢复 checkpoint，并实际演练节点失败后的恢复时间。
\end{enumerate}

\subsection{拓展阅读}

\begin{itemize}
  \item Hugging Face，\href{https://huggingface.co/spaces/nanotron/ultrascale-playbook}{Ultra-Scale Playbook}：本课的主要系统图与决策框架。
  \item Rajbhandari 等，\href{https://arxiv.org/abs/1910.02054}{ZeRO}；PyTorch，\href{https://docs.pytorch.org/docs/stable/distributed.fsdp.fully_shard.html}{FSDP2 fully\_shard documentation}。
  \item Liu 等，\href{https://arxiv.org/abs/2310.01889}{Ring Attention with Blockwise Transformers}。
  \item DeepSeek-AI，\href{https://arxiv.org/abs/2412.19437}{DeepSeek-V3 Technical Report} 与 \href{https://github.com/deepseek-ai/DeepEP}{DeepEP}。
  \item \href{https://github.com/pytorch/torchtitan}{TorchTitan}、\href{https://github.com/huggingface/nanotron}{Nanotron}、\href{https://github.com/NVIDIA/Megatron-LM}{Megatron-LM}：可执行的多维并行实现。
  \item Hugging Face，\href{https://huggingface.co/spaces/HuggingFaceTB/smol-training-playbook}{Smol Training Playbook}；\href{https://jax-ml.github.io/scaling-book/tpus/}{JAX Scaling Book}：训练基础设施与 TPU 视角。
\end{itemize}

\end{document}
